
\Needspace{12\baselineskip}
\subsubsection*{Naučnoistraživački rad -- \ClanPet}

\textbf{Naslov rada:}
\emph{Credit Card Fraud Detection Using Improved Deep Learning Models}.\\

\textbf{Autori:}
Sumaya S. Sulaiman, Ibraheem Nadher i Sarab M. Hameed.\\

\textbf{Izvor:}
\emph{Computers, Materials \& Continua}, vol. 78, br. 1, str. 1049--1069,
2024. \cite{sulaiman2024fraud}

\begin{table}[H]
\centering
\caption{Sažeti pregled naučnoistraživačkog rada Sulaiman i saradnika (2024)}
\label{tab:sulaiman-pregled}
\small
\renewcommand{\arraystretch}{1.22}
\begin{tabularx}{\linewidth}{@{}p{0.25\linewidth}Y@{}}
\toprule
\textbf{Stavka} & \textbf{Opis} \\
\midrule

Problem &
Detekcija prevara kreditnim karticama pomoću modela dubokog učenja
i ispitivanje uticaja podešavanja hiperparametara i metoda
uravnotežavanja na njihove performanse. \\

Skup podataka &
Evropski skup transakcija kreditnim karticama (Kaggle, septembar 2013.):
284\,807 transakcija, od kojih su 492 prevarne (0,172\,\%). \\

Korištene metode &
Autoenkoder (AE), konvolucijska neuronska mreža (CNN) i LSTM mreža;
podešavanje hiperparametara metodama \emph{random search} i Bayesovom
optimizacijom (cilj: F1-mjera); uravnotežavanje metodama RUS, SMOTE
i ADASYN; stratifikovana unakrsna validacija. \\

Metrike &
Tačnost, stopa detekcije (\emph{detection rate}, tj. recall) i AUC;
F1-mjera kao funkcija cilja pri podešavanju hiperparametara. \\

Najbolji rezultat &
\textbf{LSTM + SMOTE} (\emph{random search}): tačnost \textbf{99,2\,\%},
stopa detekcije \textbf{93,3\,\%}, AUC \textbf{96,3\,\%}. \\

\bottomrule
\end{tabularx}
\end{table}

\textbf{Cilj i istraživački problem.}
Autori polaze od toga da uspješnost modela dubokog učenja u velikoj
mjeri zavisi od izbora hiperparametara. Cilj je ispitati kako
sistematsko podešavanje (broj neurona, \emph{dropout}, aktivacijska
funkcija, veličina \emph{batch}-a, optimizator, funkcija gubitka i
stopa učenja) u kombinaciji s različitim metodama uravnotežavanja
utiče na detekciju prevara.

\textbf{Skup podataka i priprema.}
Korišten je javni evropski skup sa 30 numeričkih atributa i izrazitom
nebalansiranošću klasa. Podaci su podijeljeni na skup za treniranje,
validaciju i testiranje stratifikovanom unakrsnom validacijom, a
uravnotežavanje (RUS, SMOTE ili ADASYN) primijenjeno je \emph{samo
na skupu za treniranje}. Za svaki model izvedeno je osam pokretanja
(četiri slučajnom pretragom i četiri Bayesovom optimizacijom), uz
četiri scenarija: bez uravnotežavanja, RUS, SMOTE i ADASYN.

\textbf{Poređenje eksperimentalnih rezultata.}
Tabela~\ref{tab:sulaiman-rezultati} prikazuje dio rezultata iz Tabele 3
rada (podešavanje hiperparametara slučajnom pretragom).

\begin{table}[H]
\centering
\caption{Rezultati modela dubokog učenja iz rada Sulaiman i saradnika (u \%)}
\label{tab:sulaiman-rezultati}
\small
\setlength{\tabcolsep}{4pt}
\renewcommand{\arraystretch}{1.18}
\begin{tabularx}{\linewidth}{@{}Y c c c@{}}
\toprule
\textbf{Model -- uravnotežavanje} & \textbf{Acc.} & \textbf{Stopa detekcije} & \textbf{AUC} \\
\midrule
AE -- bez uravnotežavanja & 95,0 & 90,0 & 92,7 \\
AE -- SMOTE & 96,1 & 89,7 & 92,9 \\
CNN -- bez uravnotežavanja & 99,7 & 83,0 & 91,4 \\
CNN -- SMOTE & 99,8 & 89,7 & 94,7 \\
LSTM -- bez uravnotežavanja & 99,9 & 80,8 & 90,4 \\
LSTM -- RUS & 98,4 & 88,9 & 93,6 \\
\textbf{LSTM -- SMOTE} & \textbf{99,2} & \textbf{93,3} & \textbf{96,3} \\
LSTM -- ADASYN & 99,2 & 92,6 & 95,9 \\
\bottomrule
\end{tabularx}
\end{table}

\textbf{Tumačenje rezultata.}
Uravnotežavanje skupa za treniranje značajno povećava stopu detekcije:
kod LSTM mreže ona raste s 80,8\,\% (bez uravnotežavanja) na 93,3\,\%
uz SMOTE, uz mali pad tačnosti (s 99,9\,\% na 99,2\,\%). Metode koje
generišu sintetičke primjere (SMOTE, ADASYN) bile su uspješnije od
RUS-a, koji odbacuje dio podataka većinske klase. Autori LSTM i CNN
modelima pripisuju prednost jer mogu naučiti sekvencijalne obrasce
transakcija. Rad ne prikazuje preciznost, pa nije moguće procijeniti
broj lažnih alarma koji prati visoku stopu detekcije.

\textbf{Zaključak i značaj za istraživanje.}
Rad pokazuje da podešavanje hiperparametara i oversampling skupa
za treniranje mogu značajno povećati osjetljivost modela, što je
direktno povezano sa zadacima 2.2 našeg projekta (\emph{tuning} i
uravnotežavanje klasa). Istovremeno, potvrđuje važnost prikaza
preciznosti uz recall: bez preciznosti visoka stopa detekcije ne
govori koliko je model upotrebljiv u praksi.
